\chapter{Active Learning}
\label{ch:active-learning}

A potential that explores a reaction network must be accurate on exactly the
configurations the exploration visits: strained geometries, partially formed
bonds, transition-state regions.  A foundation model is not, and the failure
is systematic rather than random -- on the development benchmark the barrier
heights carried a fitted slope of 0.588, a strong compression that a
post-hoc linear correction cannot repair across systems.  The protocol used
here therefore turns explored paths into labelled data and labelled data into
a fine-tuned potential, in a five-step loop: sample, label, split, fine-tune,
verify.  Each step has a stand-alone tool, listed in
Chapter~\ref{ch:scripts}.

Related strategies appear in the literature: iteratively trained neural
network potentials for reaction-path search \cite{staub2025}, lifelong
machine-learning potentials \cite{eckhoff2025parrot}, and analyses of where
NNPs fail for kinetics predictions \cite{staub2023}.  The protocol below is
the one implemented and measured in this package.

\section{The five-step loop}
\label{sec:al-loop}

\begin{table}[htbp]
\centering
\caption{The active-learning loop and the tool for each step.  Sampling and
labelling are the two halves of one script; fine-tuning itself is performed
by the external trainer.}
\label{tab:al-steps}
\begin{tabularx}{\textwidth}{@{}llX@{}}
\toprule
Step & Tool & What happens \\
\midrule
1.\ Sample & \script{build_al_trainset} & configurations along the MEP, via \code{spline.evaluate(t)} \\
2.\ Label & \script{build_al_trainset} & ORCA L2 energies and forces, in parallel \\
3.\ Split & \script{resplit_trainset} & composition-stratified, reaction-grouped \\
4.\ Fine-tune & (external trainer) & MACE fine-tuning on the labelled frames \\
5.\ Verify & \script{validate_finetune} & barrier MAE and bias decomposition on held-out reactions \\
\bottomrule
\end{tabularx}
\end{table}

\subsection{Step 1: sampling the paths}
\label{sec:al-sample}

The sampling unit is the minimum-energy path the search has already refined.
Each elementary step in the database carries a spline fit of that path, and
\code{spline.evaluate(t)} returns a complete geometry at a fraction $t$ of the
reaction coordinate: $t=0$ is the reactant, $t$ at the spline's
\code{ts_position} is the transition state.  Sampling through the spline call,
rather than re-parsing the stored coordinate arrays, keeps the frames
consistent with the path.

A uniform grid wastes its points: at fifteen points only one or two land near
the transition state, which is exactly where the energy is most sensitive to
the barrier and the forces are largest.  \script{build_al_trainset} therefore
places half of the requested points inside a window of $\pm0.25$ around the
transition state and distributes the rest over the full path
(\code{sample_ts_weights}); the \code{--uniform} flag restores the uniform
grid.

\subsection{Step 2: labelling}
\label{sec:al-label}

Every sampled frame is labelled with ORCA at level L2
(r2SCAN-D4/def2-SVP \cite{furness2020r2scan,caldeweyher2019d4,neese2022orca}).
The labelling runs in worker processes and writes one frame per configuration
in extended-xyz format, with energies and Cartesian forces in eV and \AA, the
element list, and the bookkeeping fields (\code{es_id}, $t$) that the later
steps use to group and stratify.  Frames whose reference calculation fails are
skipped and counted rather than silently replaced.

\subsection{Step 3: splitting}
\label{sec:al-split}

The split must be stratified by composition and grouped by reaction.  Two
failure modes motivate this.  An initial split that simply took the tail of
the elementary-step order produced a training set entirely of C$_4$H$_6$ and a
validation set entirely of C$_6$H$_{10}$; the trainer scored the validation
set with the training set's atomic energies and reported an energy RMSE of
29~eV/atom.  And a random split of frames would place consecutive frames of
one path -- nearly identical geometries -- on both sides of the split.
\script{resplit_trainset} merges the available frame batches and re-splits
them: frames are grouped by elementary step (\code{es_id}), and the steps are
divided 8:2 per composition so that both sides share the same composition
distribution.

\subsection{Step 4: fine-tuning}
\label{sec:al-finetune}

Fine-tuning uses MACE \cite{batatia2022mace} on the labelled frames, with the
held-out reactions as the validation set -- never random frames of trained
reactions.  Two practical points are worth recording, because both failure
modes are opaque:

\begin{itemize}
  \item the frames carry the keys \code{energy} and \code{forces}; if the
        trainer's default key names differ, it must be told to read these;
  \item the model's element coverage follows the training data, so the
        training and validation sets must cover every element of the target
        chemistry.  A potential fine-tuned on C/H data alone will refuse
        configurations that contain oxygen.
\end{itemize}

\subsection{Step 5: verification}
\label{sec:al-verify}

Verification re-scores the held-out reactions and reports the mean absolute
error of the barrier heights together with its decomposition: by composition,
by reaction type, and by system size.  For bimolecular elementary steps, where
the path cannot be rebuilt by concatenating the reactant fragments, the
barrier is evaluated from the spline endpoints -- reactant at $t=0$,
transition state at the spline's \code{ts_position} -- so that the reference
and every model see the same two geometries
(\script{eval_barrier_spline}).  \script{validate_finetune} performs the
comparison between the fine-tuned model, the foundation model and the ORCA
reference on the held-out frames.

\section{Three lessons}
\label{sec:al-lessons}

\subsection{Training composition determines generality}
A single-composition training set produces a specialist: excellent inside its
composition, worse than the untuned foundation model outside it.  Extending
the training composition removes that specialisation.

\subsection{Reaction-type coverage is a separate axis}
A model fine-tuned on single-molecule rearrangements can still degrade on
bimolecular reactions even after composition extension.  Reaction-type
coverage is independent of composition coverage, and the out-of-distribution
check belongs to the release protocol rather than to the optional extras.

\subsection{Accuracy is not uniform across system size}
Report errors in size bins.  An aggregate mean hides systematically weak bins,
and a potential that is good on average can be unusable for the size class
that matters to a given study.

\section{Measured record}
\label{sec:al-record}

On the development systems (C/H/N/O chemistry, ORCA L2 as the reference,
barriers in kcal/mol):

\begin{itemize}
  \item the fine-tuned potential, on held-out reactions with
        cross-composition training, reaches a barrier MAE of 2.92, against
        6.93 for the strongest untuned foundation model on the same benchmark;
  \item the committee disagreement tracks the true force error with a rank
        correlation of 0.847 at calibration, and the fallback trigger fires
        preferentially in transition-state regions, where the disagreement
        rises by several times over its equilibrium value; this second number
        belongs to the uncertainty layer, described in the companion chapter
        on uncertainty quantification.
\end{itemize}

\begin{readbox}[How to read these numbers]
They characterise the protocol on the tested systems; they are not
transferable accuracy guarantees for other chemistries.  Reproduce them on
your own reference set before quoting them for a new system, and quote the
error decomposition alongside the aggregate.
\end{readbox}
